\documentclass[11pt,a4paper]{article}
% preamble.tex — shared setup for human-verification documents
% Compile each document with:  tectonic <File>.tex

\usepackage[T1]{fontenc}
\usepackage[utf8]{inputenc}
\usepackage{lmodern}
\usepackage{amssymb,amsmath,amsthm}
\usepackage{mathtools}
\usepackage[dvipsnames]{xcolor}
\usepackage{listings}
\usepackage{hyperref}
\usepackage[margin=2.5cm]{geometry}
\usepackage{parskip}

% ── Colours ──────────────────────────────────────────────────────────
\definecolor{leanBg}{HTML}{F7F7F0}
\definecolor{haskellBg}{HTML}{F0F4F7}
\definecolor{keyword}{HTML}{7B3294}
\definecolor{comment}{HTML}{5A8C32}
\definecolor{noteFg}{HTML}{B35900}

% ── Lean 4 listings language ─────────────────────────────────────────
% NOTE: Lean listing blocks use ASCII approximations of unicode symbols
% (e.g. -> for arrow, o. for compose).  The literate key substitutes
% these with proper math-mode renderings in the PDF.
\lstdefinelanguage{Lean4}{
  morekeywords={def,theorem,structure,where,namespace,end,import,fun,
                let,have,show,by,exact,intro,cases,rfl,subst,obtain,
                funext,simp,rw,apply,variable,Type,Prop,instance,
                section,open,class,extends,deriving,sorry,noncomputable},
  sensitive=true,
  morecomment=[l]{--},
  morecomment=[n]{/-}{-/},
  morestring=[b]",
}

% ── Lean listing style ───────────────────────────────────────────────
\lstdefinestyle{lean}{
  language=Lean4,
  backgroundcolor=\color{leanBg},
  basicstyle=\small\ttfamily,
  keywordstyle=\color{keyword}\bfseries,
  commentstyle=\color{comment}\itshape,
  breaklines=true,
  frame=single,
  framerule=0.4pt,
  xleftmargin=4pt,
  xrightmargin=4pt,
  aboveskip=6pt,
  belowskip=4pt,
  literate=
    {->}{{$\to$}}2
    {=>}{{$\Rightarrow$}}2
    {<|}{{$\langle$}}1
    {|>}{{$\rangle$}}1
    {==>}{{$\Longrightarrow$}}3
    { o. }{{$\;\circ\;$}}3
    {forall}{{$\forall$}}1
    {:=}{{$\,:=\,$}}2
    {!=}{{$\neq$}}2
    {<=}{{$\leq$}}2
    {>=}{{$\geq$}}2,
}

% ── Haskell listing style ───────────────────────────────────────────
\lstdefinestyle{haskell}{
  language=Haskell,
  backgroundcolor=\color{haskellBg},
  basicstyle=\small\ttfamily,
  keywordstyle=\color{keyword}\bfseries,
  commentstyle=\color{comment}\itshape,
  breaklines=true,
  frame=single,
  framerule=0.4pt,
  xleftmargin=4pt,
  xrightmargin=4pt,
  aboveskip=6pt,
  belowskip=4pt,
  literate=
    {->}{{$\to$}}2
    {=>}{{$\Rightarrow$}}2,
}

% ── Theorem environments ────────────────────────────────────────────
\theoremstyle{definition}
\newtheorem{defn}{Definition}[section]
\newtheorem{thm}[defn]{Theorem}
\newtheorem{prop}[defn]{Proposition}

\theoremstyle{remark}
\newtheorem*{fidelitynote}{Fidelity note}

% ── Convenience macros ───────────────────────────────────────────────
\newcommand{\lean}[1]{\texttt{#1}}      % inline Lean identifier
\newcommand{\hask}[1]{\texttt{#1}}      % inline Haskell identifier
\newcommand{\Ext}{\mathrm{Ext}}
\newcommand{\Pos}{\mathrm{Pos}}
\newcommand{\Shape}{\mathrm{Shape}}


\title{Human Verification: \texttt{TrajectoryComposition}\\[4pt]
  \large Composing Trajectory Evaluation with Trust Evaluation}
\author{Machine-generated from \texttt{lean/TrajectoryComposition.lean} for review by Neil Ghani}
\date{July 2026}

\begin{document}
\maketitle

\begin{abstract}
This document presents every definition and theorem \emph{statement} from the
Lean~4 module \lean{TrajectoryComposition} in three parallel forms: the original
Lean, a readable Haskell rendering, and precise mathematical English.  Proofs
are omitted --- they are machine-checked by Lean.

The module defines a grade-level map $\psi : \mathtt{TLevel} \to
\mathtt{TrustLevel}$ from the five-level trajectory chain to the seven-level
deployed trust chain, verifies that $\psi$ is a chain morphism (monotone,
bottom-preserving), and applies the lax transport theorem
(\lean{tau\_transport} from \lean{Trust.lean}) to concrete trajectories.
\end{abstract}

\tableofcontents
\bigskip

%===================================================================
\section{Overview}
%===================================================================

The module bridges two trust chains.  On one side is the five-level trajectory
chain \lean{TLevel} (defined in \lean{Trajectory.lean}): $\mathtt{unsupported}
< \mathtt{shallow} < \mathtt{independent} < \mathtt{crossVerified} <
\mathtt{mechChecked}$.  On the other is the seven-level deployed agent chain
\lean{TrustLevel} used in production registries: $\mathtt{unchecked} <
\mathtt{speculative} < \mathtt{computed} < \mathtt{proved} <
\mathtt{peerReviewed} < \mathtt{published} < \mathtt{leanVerified}$.

The map $\psi$ translates: $\mathtt{unsupported} \mapsto \mathtt{unchecked}$,
$\mathtt{shallow} \mapsto \mathtt{computed}$, $\mathtt{independent} \mapsto
\mathtt{proved}$, $\mathtt{crossVerified} \mapsto \mathtt{peerReviewed}$,
$\mathtt{mechChecked} \mapsto \mathtt{leanVerified}$.

%===================================================================
\section{The trust chain}
%===================================================================

\subsection{Inductive type: \lean{TrustLevel} (local copy)}

\begin{defn}[Deployed trust levels]
A seven-element chain mirroring \lean{TrustBoundary.TrustLevel}:
\[
  \mathtt{unchecked} < \mathtt{speculative} < \mathtt{computed} <
  \mathtt{proved} < \mathtt{peerReviewed} < \mathtt{published} <
  \mathtt{leanVerified}.
\]
\end{defn}

\paragraph{Lean.}
\begin{lstlisting}[style=lean]
inductive TrustLevel where
  | unchecked
  | speculative
  | computed
  | proved
  | peerReviewed
  | published
  | leanVerified
deriving DecidableEq, Repr
\end{lstlisting}

\paragraph{Haskell.}
\begin{lstlisting}[style=haskell]
data TrustLevel
  = Unchecked | Speculative | Computed | Proved
  | PeerReviewed | Published | LeanVerified
  deriving (Eq, Ord, Show)
\end{lstlisting}

\begin{fidelitynote}
This type is defined locally in \lean{Trajectory} namespace, mirroring the
one in \lean{TrustBoundary.lean}.  The \lean{Chain TrustLevel} instance here
is also proved locally.
\end{fidelitynote}


\subsection{Definition: \lean{TrustLevel.le} and \lean{Chain TrustLevel}}

\begin{defn}[Order on \lean{TrustLevel}]
The order is defined case-by-case, establishing the chain
$\mathtt{unchecked} \le \cdots \le \mathtt{leanVerified}$.  The
\lean{Chain TrustLevel} instance has $\bot = \mathtt{unchecked}$.
\end{defn}

\paragraph{Lean.}
\begin{lstlisting}[style=lean]
def TrustLevel.le : TrustLevel -> TrustLevel -> Prop
  | .unchecked, _ => True
  | .speculative, .unchecked => False
  | .speculative, _ => True
  -- ... (remaining cases follow the chain order)
  | .leanVerified, .leanVerified => True
  | .leanVerified, _ => False

instance : Chain TrustLevel where
  bot := .unchecked
  -- reflexivity, transitivity, antisymmetry, totality, bot_le, dec all proved
\end{lstlisting}

\paragraph{Haskell.}
\begin{lstlisting}[style=haskell]
-- Derived Ord suffices; bot = Unchecked.
\end{lstlisting}

%===================================================================
\section{The trajectory-trust morphism}
%===================================================================

\subsection{Definition: \lean{psi}}

\begin{defn}[Grade translation map]
The function $\psi : \mathtt{TLevel} \to \mathtt{TrustLevel}$ translates
trajectory grades to deployed trust grades:
\begin{align*}
  \psi(\mathtt{unsupported})  &= \mathtt{unchecked}, \\
  \psi(\mathtt{shallow})       &= \mathtt{computed}, \\
  \psi(\mathtt{independent})   &= \mathtt{proved}, \\
  \psi(\mathtt{crossVerified}) &= \mathtt{peerReviewed}, \\
  \psi(\mathtt{mechChecked})   &= \mathtt{leanVerified}.
\end{align*}
\end{defn}

\paragraph{Lean.}
\begin{lstlisting}[style=lean]
def psi : TLevel -> TrustLevel
  | .unsupported  => .unchecked
  | .shallow       => .computed
  | .independent   => .proved
  | .crossVerified => .peerReviewed
  | .mechChecked   => .leanVerified
\end{lstlisting}

\paragraph{Haskell.}
\begin{lstlisting}[style=haskell]
psi :: TLevel -> TrustLevel
psi Unsupported  = Unchecked
psi Shallow       = Computed
psi Independent   = Proved
psi CrossVerified = PeerReviewed
psi MechChecked   = LeanVerified
\end{lstlisting}


\subsection{Theorem: \lean{psi\_bot}}

\begin{thm}[$\psi$ preserves bottom]
\[
  \psi(\bot_{\mathtt{TLevel}}) = \bot_{\mathtt{TrustLevel}},
  \quad\text{i.e.,}\quad
  \psi(\mathtt{unsupported}) = \mathtt{unchecked}.
\]
This is the strictness-at-bottom condition required for $\psi$ to be a valid
trust-system morphism (see \lean{Mor.phi\_bot} in \lean{Trust.lean}).
\end{thm}

\paragraph{Lean.}
\begin{lstlisting}[style=lean]
theorem psi_bot :
    psi (Chain.bot : TLevel) = (Chain.bot : TrustLevel)
\end{lstlisting}

\paragraph{Haskell.}
\begin{lstlisting}[style=haskell]
-- psi Unsupported == Unchecked   -- True by definition
\end{lstlisting}


\subsection{Theorem: \lean{psi\_mono}}

\begin{thm}[$\psi$ is monotone]
For all $a, b : \mathtt{TLevel}$,
\[
  a \le b \implies \psi(a) \le \psi(b).
\]
Together with \lean{psi\_bot}, this makes $\psi$ a valid \lean{phi}
component of a \lean{Mor} in the sense of \lean{Trust.lean}.
\end{thm}

\paragraph{Lean.}
\begin{lstlisting}[style=lean]
theorem psi_mono : forall {a b : TLevel}, a <= b -> psi a <= psi b
\end{lstlisting}

\paragraph{Haskell.}
\begin{lstlisting}[style=haskell]
-- psi is order-preserving: a <= b ==> psi a <= psi b
-- Follows from the chain structure and the explicit case table of psi.
\end{lstlisting}

%===================================================================
\section{Concrete transport results}
%===================================================================

\subsection{Theorem: \lean{psi\_circularTraj\_bot}}

\begin{thm}[Circular trajectory maps to unchecked under $\psi$]
For every event $e : \mathtt{Evt}$,
\[
  \psi\bigl(\tau(\mathtt{circularTraj},\, e)\bigr) = \mathtt{unchecked}.
\]
This follows immediately from \lean{circularTraj\_tau\_bot} and
\lean{psi\_bot}.
\end{thm}

\paragraph{Lean.}
\begin{lstlisting}[style=lean]
theorem psi_circularTraj_bot : forall (e : Evt),
    psi (tau circularTraj circularTraj_wf e) = TrustLevel.unchecked
\end{lstlisting}

\paragraph{Haskell.}
\begin{lstlisting}[style=haskell]
-- psi (tau circularTraj e) == Unchecked  for all e
\end{lstlisting}


\subsection{Theorem: \lean{psi\_honest\_e4}}

\begin{thm}[Honest trajectory's verify event maps to proved]
\[
  \psi\bigl(\tau(\mathtt{honestTraj},\, e_4)\bigr) = \mathtt{proved}.
\]
Recall from \lean{Trajectory.lean} that $\tau(\mathtt{honestTraj},\, e_4) =
\mathtt{independent}$, and $\psi(\mathtt{independent}) = \mathtt{proved}$.
\end{thm}

\paragraph{Lean.}
\begin{lstlisting}[style=lean]
theorem psi_honest_e4 :
    psi (tau honestTraj honestTraj_wf Evt.e4) = TrustLevel.proved
\end{lstlisting}

\paragraph{Haskell.}
\begin{lstlisting}[style=haskell]
-- psi (tau honestTraj E4) == Proved
\end{lstlisting}


\subsection{Theorem: \lean{psi\_honest\_e3}}

\begin{thm}[Honest trajectory's claim event maps to computed]
\[
  \psi\bigl(\tau(\mathtt{honestTraj},\, e_3)\bigr) = \mathtt{computed}.
\]
Recall that $\tau(\mathtt{honestTraj},\, e_3) = \mathtt{shallow}$, and
$\psi(\mathtt{shallow}) = \mathtt{computed}$.
\end{thm}

\paragraph{Lean.}
\begin{lstlisting}[style=lean]
theorem psi_honest_e3 :
    psi (tau honestTraj honestTraj_wf Evt.e3) = TrustLevel.computed
\end{lstlisting}

\paragraph{Haskell.}
\begin{lstlisting}[style=haskell]
-- psi (tau honestTraj E3) == Computed
\end{lstlisting}

%===================================================================
\section{Translation conventions}
%===================================================================

\begin{enumerate}
\item \textbf{Proof bodies omitted.}  Lean machine-checks all proofs.

\item \textbf{Local type duplication.}  \lean{TrustLevel} is defined
  independently here (mirroring \lean{TrustBoundary.TrustLevel}) to keep the
  file self-contained and avoid namespace friction.  A future refactor could
  import from \lean{TrustBoundary.lean} instead.

\item \textbf{The morphism condition.}  \lean{psi\_bot} and \lean{psi\_mono}
  are exactly the \lean{phi\_bot} and \lean{phi\_mono} fields required to
  extend $\psi$ to a \lean{Mor} between trust systems.  The transport theorem
  (\lean{tau\_transport}) then applies: importing along such a morphism never
  claims more than $\psi$ of the source's canonical trust.

\item \textbf{Concrete transport.}  Rather than constructing a full
  \lean{Mor} record (which would require matching step-sets and premise maps
  between the trajectory and trust systems), the module applies results from
  \lean{Trajectory.lean} directly, composing with $\psi$ pointwise.
\end{enumerate}

\end{document}
